\documentclass[sigconf]{acmart}

\settopmatter{printacmref=false}
\renewcommand\footnotetextcopyrightpermission[1]{}
\pagestyle{plain}

\usepackage{booktabs}
\usepackage{array}

% Copy the exported appendix figures into a local `figures/` directory and use
% the simplified filenames referenced below:
%
%   appendix_recall_vs_qps.pdf
%   appendix_recall_vs_index_overhead.pdf
%   appendix_init_recall_vs_avg_computations.pdf
%   appendix_query_ram_vs_index_overhead.pdf
%   appendix_serving_ram_vs_index_overhead.pdf
%   appendix_computation_range_bars.pdf
%   appendix_ablation_recall_vs_avg_computations.pdf
%   appendix_ablation_recall_vs_index_overhead.pdf

\graphicspath{{figures/}}

\begin{document}

\title{Supplementary Material: Uncertainty-Aware Learned Routing for Partitioned Vector Search}
\author{Anonymous Authors}
\maketitle

\section{Additional Tradeoff Views}

\noindent This appendix adds views that are informative but too detailed for the
main paper. The figures in this section refine the competitiveness story along
three complementary axes: throughput, structural cost, and workload-normalized
initialization behavior. Together they show that the strongest AdaptIVF regime
is obtained by changing query-time effort on a compact backbone, not by moving
to a heavier serving structure.

\begin{figure*}[!htbp]
  \centering
  \includegraphics[width=\linewidth]{appendix_recall_vs_qps.pdf}
  \caption{Recall@10 versus QPS across the full competitiveness study. This view isolates the quality--throughput frontier directly: classical baselines remain the raw speed anchors, while the stronger AdaptIVF settings trade additional online work for higher recall rather than for larger index structures.}
  \Description{Faceted scatter plot with four dataset panels: GloVe, SIFT, GIST, and Deep1M. The x-axis is QPS and the y-axis is Recall@10. HNSW and IVFPQ appear at high-throughput regions, while learned methods occupy lower-throughput zones. AdaptIVF-80 and AdaptIVF+PQ-80 shift upward in recall compared with their 10-cap counterparts, but also move toward lower throughput, visualizing a recall-throughput tradeoff within the same method family.}
  \label{fig:appendix:recall_qps}
\end{figure*}

\begin{figure*}[!htbp]
  \centering
  \includegraphics[width=\linewidth]{appendix_recall_vs_index_overhead.pdf}
  \caption{Recall@10 versus index overhead. This figure makes the structural split among the learned methods explicit: LIRA reaches strong recall with a much heavier index, whereas AdaptIVF retains a substantially smaller footprint while recovering high-cap recall.}
  \Description{Faceted scatter plot across GloVe, SIFT, GIST, and Deep1M. The x-axis is index overhead in MB and the y-axis is Recall@10. LIRA sits far to the right with high overhead, while AdaptIVF and AdaptIVF+PQ remain much farther left. Increasing probe cap from 10 to 80 for AdaptIVF variants raises recall with nearly unchanged horizontal position, indicating recall gains without major structural growth.}
  \label{fig:appendix:recall_index}
\end{figure*}

\begin{figure*}[!htbp]
  \centering
  \includegraphics[width=\linewidth]{appendix_init_recall_vs_avg_computations.pdf}
  \caption{Initialization study under the workload axis. Expressing the sweep against average computations rather than probe depth preserves the same qualitative split as the main initialization figure: geometry-aware initialization is decisive on SIFT, while BLISS remains competitive or stronger on GIST and Deep1M.}
  \Description{Faceted plot with four dataset panels showing recall versus average computations for BLISS, BLISS-KMeans, and MLP-IVF. SIFT shows BLISS-KMeans and MLP-IVF above BLISS at similar or lower computation ranges. GIST and Deep1M show BLISS competitive with or above the geometry-initialized variants. GloVe shows smaller differences among methods.}
  \label{fig:appendix:init_avg_comp}
\end{figure*}

\noindent Read together, these three figures refine the main-paper account
without changing it. The first two confirm that the strongest AdaptIVF
operating regime is reached by spending more online work on the same compact
index, while the initialization workload plot shows that the choice of coarse
seed is genuinely dataset-dependent rather than uniformly favorable to
geometric seeding.

\section{Query Work Distribution}

\noindent Average computations are the main workload metric in the paper, but
they can hide substantial per-query variation. The following figure adds that
missing detail by contrasting minimum, average, and maximum computation counts
for the methods that report all three quantities. This is especially relevant
for adaptive routing, where the workload is meant to respond to query
difficulty rather than stay fixed. The compared methods here are
\textbf{BLISS}, \textbf{LIRA}, \textbf{AdaptIVF-10}, \textbf{AdaptIVF-80},
\textbf{AdaptIVF+PQ-10}, and \textbf{AdaptIVF+PQ-80}.

\begin{figure*}[!htbp]
  \centering
  \includegraphics[width=\linewidth]{appendix_computation_range_bars.pdf}
  \caption{Per-query computation range for \textbf{BLISS}, \textbf{LIRA}, \textbf{AdaptIVF-10}, \textbf{AdaptIVF-80}, \textbf{AdaptIVF+PQ-10}, and \textbf{AdaptIVF+PQ-80}. This figure complements the average-computation plots by showing that methods with similar means can have very different tails. In particular, the higher-cap AdaptIVF variants expand the upper tail substantially, which is consistent with their stronger recall regime and clarifies that the main tradeoff is concentrated in hard queries rather than in uniform work inflation.}
  \Description{Grouped bar-style comparison across datasets showing minimum, average, and maximum computation counts for selected methods. Low-cap settings have tighter ranges, while high-cap AdaptIVF variants show visibly larger maximum bars, indicating heavier worst-case work. The figure emphasizes that tail behavior differs even when averages are similar.}
  \label{fig:appendix:computation_range}
\end{figure*}

\noindent The main qualitative pattern is that low-cap configurations remain
closer to a bounded workload regime, whereas the stronger high-cap settings
derive part of their recall gain from allowing much larger worst-case search
effort. That is precisely the controller-calibration question left open by the
main results: how much of that tail can be retained for difficult queries
without paying for it as often.

\section{Memory and Serving Diagnostics}

\noindent The main paper argues that AdaptIVF should be understood as a compact
learned system, not merely as a retrieval method. The memory diagnostics below
back that claim from two sides: query-time RSS deltas capture incremental
online cost, while serving-overhead RSS measures the resident cost of loading
the required artifacts before queries arrive. The compared methods are
\textbf{BLISS}, \textbf{LIRA}, \textbf{AdaptIVF-10}, \textbf{AdaptIVF-80},
\textbf{AdaptIVF+PQ-10}, and \textbf{AdaptIVF+PQ-80}.

\begin{figure*}[!htbp]
  \centering
  \includegraphics[width=\linewidth]{appendix_query_ram_vs_index_overhead.pdf}
  \vfill
  \includegraphics[width=\linewidth]{appendix_serving_ram_vs_index_overhead.pdf}
  \caption{Memory diagnostics for \textbf{BLISS}, \textbf{LIRA}, \textbf{AdaptIVF-10}, \textbf{AdaptIVF-80}, \textbf{AdaptIVF+PQ-10}, and \textbf{AdaptIVF+PQ-80}. Top: isolated query-path RAM delta versus index overhead. Bottom: isolated post-load serving RAM versus index overhead. The same pattern appears on both views: the single-backbone AdaptIVF family remains materially lighter than LIRA, while the compressed baselines retain the smallest memory footprint overall.}
  \Description{Two stacked scatter plots. The top plot maps query-path RAM delta against index overhead. The bottom plot maps post-load serving RAM overhead against index overhead. Across both plots, LIRA appears with substantially higher memory and overhead, AdaptIVF variants cluster at lower overhead, and compressed variants occupy lower-memory regions. The paired views separate online per-query memory cost from resident serving memory cost.}
  \label{fig:appendix:memory}
\end{figure*}

\begin{table}[!htbp]
  \centering
  \scriptsize
  \setlength{\tabcolsep}{4pt}
  \caption{Learned-method memory summary. Query RAM is the isolated per-query RSS delta; Serving Overhead is the isolated post-load method overhead in MB. The table shows that AdaptIVF remains materially lighter than LIRA both at query time and after the serving artifacts are loaded.}
  \label{tab:appendix:memory}
  \begin{tabular}{llrr}
    \toprule
    Dataset & Method & Query RAM & Serving Overhead \\
    \midrule
    GloVe & BLISS & 24.1 & 507.6 \\
    GloVe & LIRA & 25.0 & 989.4 \\
    GloVe & AdaptIVF-10 & 12.1 & 498.2 \\
    GloVe & AdaptIVF-80 & 38.2 & 498.0 \\
    GloVe & AdaptIVF+PQ-10 & 10.2 & 516.8 \\
    GloVe & AdaptIVF+PQ-80 & 15.2 & 516.4 \\
    \midrule
    SIFT & BLISS & 31.5 & 540.0 \\
    SIFT & LIRA & 169.4 & 1020.9 \\
    SIFT & AdaptIVF-10 & 17.0 & 531.4 \\
    SIFT & AdaptIVF-80 & 41.0 & 531.5 \\
    SIFT & AdaptIVF+PQ-10 & 10.9 & 546.9 \\
    SIFT & AdaptIVF+PQ-80 & 12.4 & 547.0 \\
    \midrule
    GIST & BLISS & 41.7 & 3714.0 \\
    GIST & LIRA & 25.6 & 4224.9 \\
    GIST & AdaptIVF-10 & 33.6 & 3705.4 \\
    GIST & AdaptIVF-80 & 27.0 & 3704.8 \\
    GIST & AdaptIVF+PQ-10 & 9.3 & 3721.6 \\
    GIST & AdaptIVF+PQ-80 & 10.7 & 3721.6 \\
    \midrule
    Deep1M & BLISS & 45.1 & 3877.6 \\
    Deep1M & LIRA & 799.2 & 4281.1 \\
    Deep1M & AdaptIVF-10 & 30.2 & 3829.6 \\
    Deep1M & AdaptIVF-80 & 47.3 & 3829.7 \\
    Deep1M & AdaptIVF+PQ-10 & 16.1 & 3982.2 \\
    Deep1M & AdaptIVF+PQ-80 & 12.0 & 3982.1 \\
    \bottomrule
  \end{tabular}
\end{table}

\noindent The memory views are useful because they distinguish ``small on disk''
from ``light to serve.'' On both criteria, AdaptIVF remains materially closer
to BLISS than to LIRA, while still offering a more favorable storage profile
than either repeated-structure or local-index learned baselines.

\section{Detailed Tables}

\noindent The tables in this section are included for direct numerical lookup.
They are intentionally compact, but together they cover the three quantitative
surfaces most likely to be cross-checked during review: initialization at the
representative probe depth, competitiveness under recall/work, and
competitiveness under speed/overhead.

\begin{table}[!htbp]
  \centering
  \scriptsize
  \setlength{\tabcolsep}{4pt}
  \caption{Initialization summary at the representative operating point $m=10$. Each method entry reports Recall@10 and Avg. Comp. (k). The table mirrors the initialization figure numerically and shows the dataset-dependent reversal directly: geometry-aware initialization dominates on SIFT, while hash-seeded BLISS is stronger on GIST and Deep1M.}
  \label{tab:appendix:init_m10}
  \begin{tabular}{lrrrrrr}
    \toprule
    Dataset & \multicolumn{2}{c}{BLISS} & \multicolumn{2}{c}{BLISS-KMeans} & \multicolumn{2}{c}{MLP-IVF} \\
    \cmidrule(lr){2-3}\cmidrule(lr){4-5}\cmidrule(lr){6-7}
    & Recall & Avg. Comp. & Recall & Avg. Comp. & Recall & Avg. Comp. \\
    \midrule
    GloVe & 0.819 & 30.0 & 0.825 & 26.9 & 0.822 & 24.9 \\
    SIFT & 0.676 & 38.7 & 0.928 & 21.4 & 0.927 & 20.0 \\
    GIST & 0.834 & 35.0 & 0.756 & 35.4 & 0.752 & 35.8 \\
    Deep1M & 0.922 & 46.7 & 0.896 & 46.9 & 0.900 & 54.2 \\
    \bottomrule
  \end{tabular}
\end{table}

\begin{table*}[!htbp]
  \centering
  \scriptsize
  \setlength{\tabcolsep}{3pt}
  \caption{Competitiveness summary: Recall@10 and Avg. Comp. (k) at the fixed operating point used for the final run. This table provides the numeric counterpart to the main recall--work figure and makes the role of the `m80` variants explicit.}
  \label{tab:appendix:comp_recall}
  \begin{tabular}{lrrrrrrrr}
    \toprule
    Method & \multicolumn{2}{c}{GloVe} & \multicolumn{2}{c}{SIFT} & \multicolumn{2}{c}{GIST} & \multicolumn{2}{c}{Deep1M} \\
    \cmidrule(lr){2-3}\cmidrule(lr){4-5}\cmidrule(lr){6-7}\cmidrule(lr){8-9}
    & Recall & Avg. Comp. & Recall & Avg. Comp. & Recall & Avg. Comp. & Recall & Avg. Comp. \\
    \midrule
    HNSW & 0.897 & 4.9 & 0.996 & 3.0 & 0.914 & 3.8 & 0.981 & 3.7 \\
    IVF & 0.749 & 12.3 & 0.875 & 10.8 & 0.704 & 14.0 & 0.857 & 29.6 \\
    IVFPQ & 0.379 & 12.3 & 0.546 & 10.8 & 0.171 & 14.0 & 0.456 & 29.6 \\
    BLISS & 0.814 & 30.1 & 0.730 & 38.7 & 0.831 & 35.3 & 0.924 & 47.0 \\
    MLP-IVF & 0.822 & 25.1 & 0.929 & 20.4 & 0.760 & 35.9 & 0.902 & 54.1 \\
    MLP-IVFPQ & 0.758 & 25.2 & 0.928 & 20.5 & 0.591 & 36.0 & 0.888 & 54.2 \\
    LIRA & 0.875 & 0.8 & 0.972 & 0.8 & 0.916 & 0.9 & 0.932 & 0.8 \\
    AdaptIVF-10 & 0.736 & 21.4 & 0.825 & 17.5 & 0.568 & 19.0 & 0.822 & 45.6 \\
    AdaptIVF-80 & 0.913 & 53.8 & 0.956 & 31.4 & 0.921 & 104.8 & 0.954 & 91.1 \\
    AdaptIVF+PQ-10 & 0.692 & 21.4 & 0.814 & 17.5 & 0.477 & 19.1 & 0.813 & 45.7 \\
    AdaptIVF+PQ-80 & 0.801 & 54.3 & 0.954 & 31.8 & 0.661 & 103.9 & 0.934 & 90.9 \\
    \bottomrule
  \end{tabular}
\end{table*}

\begin{table*}[!htbp]
  \centering
  \scriptsize
  \setlength{\tabcolsep}{3pt}
  \caption{Competitiveness summary: QPS and Index Overhead (MB) at the same fixed operating point. This view complements Table~\ref{tab:appendix:comp_recall} by separating serving speed from structural cost.}
  \label{tab:appendix:comp_systems}
  \begin{tabular}{lrrrrrrrr}
    \toprule
    Method & \multicolumn{2}{c}{GloVe} & \multicolumn{2}{c}{SIFT} & \multicolumn{2}{c}{GIST} & \multicolumn{2}{c}{Deep1M} \\
    \cmidrule(lr){2-3}\cmidrule(lr){4-5}\cmidrule(lr){6-7}\cmidrule(lr){8-9}
    & QPS & Overhead & QPS & Overhead & QPS & Overhead & QPS & Overhead \\
    \midrule
    HNSW & 33648.7 & 758.6 & 46930.3 & 747.8 & 8323.0 & 3921.6 & 59596.6 & 6251.1 \\
    IVF & 40545.6 & 460.9 & 54114.6 & 496.4 & 5711.7 & 3673.5 & 27520.6 & 3736.2 \\
    IVFPQ & 128205.7 & 27.6 & 155753.6 & 23.5 & 77185.4 & 27.6 & 101208.8 & 230.3 \\
    BLISS & 23.5 & 45.0 & 19.8 & 39.6 & 13.5 & 46.1 & 37.7 & 337.8 \\
    MLP-IVF & 48.3 & 45.0 & 81.9 & 39.6 & 13.7 & 46.1 & 34.0 & 337.8 \\
    MLP-IVFPQ & 7.6 & 63.1 & 19.7 & 55.0 & 7.0 & 62.3 & 6.2 & 490.3 \\
    LIRA & 279.3 & 1531.3 & 457.3 & 1509.4 & 133.7 & 7857.0 & 324.6 & 12628.9 \\
    AdaptIVF-10 & 62.6 & 19.3 & 100.8 & 15.8 & 19.4 & 17.9 & 38.5 & 150.1 \\
    AdaptIVF-80 & 47.9 & 19.3 & 68.6 & 15.8 & 9.3 & 17.8 & 29.1 & 150.2 \\
    AdaptIVF+PQ-10 & 7.5 & 37.5 & 19.8 & 31.2 & 8.9 & 34.0 & 5.6 & 302.8 \\
    AdaptIVF+PQ-80 & 5.7 & 37.5 & 16.6 & 31.2 & 5.1 & 34.0 & 4.0 & 302.8 \\
    \bottomrule
  \end{tabular}
\end{table*}

\section{Ablation Diagnostics}

\noindent The main paper keeps the ablation compact because the central
question is architectural rather than exhaustive. The two figures below expand
that view and make the current limitation of the dynamic controller easier to
see: the backbone and selective-assignment design are already stable, whereas
the low-cap adaptive controller is not yet the strongest average-case operating
point. The compared variants are \textbf{AdaptIVF-Static},
\textbf{AdaptIVF-10}, \textbf{AdaptIVF-A4}, \textbf{AdaptIVF-Static+PQ},
\textbf{AdaptIVF+PQ-10}, and \textbf{AdaptIVF-A4+PQ}. Here \textbf{A4}
denotes the larger write-time replication setting $r_{\max}=4$, whereas the
default AdaptIVF variants use $r_{\max}=3$; this isolates whether extra
selective duplication alone can recover recall without changing the backbone or
the query-time uncertainty controller.

\begin{figure*}[!htbp]
  \centering
  \includegraphics[width=\linewidth]{appendix_ablation_recall_vs_avg_computations.pdf}
  \caption{Ablation recall versus average computations for \textbf{AdaptIVF-Static}, \textbf{AdaptIVF-10}, \textbf{AdaptIVF-A4}, \textbf{AdaptIVF-Static+PQ}, \textbf{AdaptIVF+PQ-10}, and \textbf{AdaptIVF-A4+PQ}. The static variant remains strongest on most datasets at the current low-cap operating point, while the larger-replication A4 setting usually increases work without consistently closing the gap.}
  \Description{Faceted ablation comparison plotting Recall@10 against average computations for ablation variants, including static and adaptive configurations with and without PQ. In most panels, static variants are at least as high in recall near similar computation levels, while larger-replication settings tend to move rightward with limited recall improvement.}
  \label{fig:appendix:ablation_avg}
\end{figure*}

\begin{figure*}[!htbp]
  \centering
  \includegraphics[width=\linewidth]{appendix_ablation_recall_vs_index_overhead.pdf}
  \caption{Ablation recall versus index overhead for \textbf{AdaptIVF-Static}, \textbf{AdaptIVF-10}, \textbf{AdaptIVF-A4}, \textbf{AdaptIVF-Static+PQ}, \textbf{AdaptIVF+PQ-10}, and \textbf{AdaptIVF-A4+PQ}. The same ordering persists when the ablations are viewed against structural cost: the main unresolved issue is controller calibration, not lack of index capacity.}
  \Description{Faceted ablation plot with index overhead on the x-axis and Recall@10 on the y-axis. Variants with similar overhead show different recall outcomes, and static low-cap variants frequently remain competitive or stronger. The figure indicates that current performance gaps are tied more to controller behavior than to insufficient index size.}
  \label{fig:appendix:ablation_overhead}
\end{figure*}

\noindent These ablation views sharpen the interpretation of the main table.
The current query-time entropy controller is not yet the strongest low-cap
average-case configuration, but the underlying compact backbone and
selective-assignment design are already well supported. The next improvement
target is therefore controller calibration, not a more elaborate index family.

\end{document}
